\documentclass[11pt]{article}
\usepackage[margin=0.9in]{geometry}
\usepackage{amsmath,amssymb,amsthm,booktabs,mathrsfs}
\usepackage[T1]{fontenc}
\usepackage{lmodern,microtype}
\usepackage[colorlinks=true,linkcolor=blue,urlcolor=blue]{hyperref}
\newtheorem{theorem}{Theorem}[section]
\newtheorem{lemma}[theorem]{Lemma}
\newtheorem{proposition}[theorem]{Proposition}
\newtheorem{corollary}[theorem]{Corollary}
\theoremstyle{remark}
\newtheorem{remark}[theorem]{Remark}
\newcommand{\norm}[1]{\left\|#1\right\|}
\title{Toward a four-dimensional Yang--Mills construction\\
\large Notebook 18: integrating overlapping cutoff interactions}
\author{Research working document}
\date{17 October 2026}
\begin{document}
\maketitle
\begin{abstract}
We extend the Gaussian cutoff expansion to interactions whose
supports overlap and can have arbitrarily many cells. An explicit
weighted local norm controls their expansion. The covariance forest
is applied after grouping overlapping factors, but its derivatives
retain their original cell endpoints. This prevents uncontrolled
powers of the grouped support size at high-degree vertices. We
obtain a nonvanishing partition function, a holomorphic logarithm
with summable activities, and a bound linear in the interaction norm
for its difference from the pure Gaussian cutoff logarithm.
A support-algebra construction also gives strictly local background
activities when the precision has a small local decomposition.
For a real baseline, a tunable relative Gaussian comparison reduces
the density cost per cell from a fixed exponential in its dimension
to an exponential with an arbitrarily small coefficient. The
theorems specify the covariance and coefficient estimates that
the Yang--Mills quantum step must supply, and a final section
identifies the corresponding original Ba\l aban integration route.
\end{abstract}

\section{The integration problem addressed here}
\label{oc18-target}

The nonlinear Wilson action and its coordinate factors generate
interactions supported on several overlapping cells. Treating each
connected collection as a single large Gaussian cell can introduce
a loss growing with both its support size and its branching degree.
The proof below keeps the original cell structure through the
covariance derivatives and sums the resulting degree factorials
exactly. The original Gaussian cell dimension stays fixed.

This is an abstract integration theorem with explicit hypotheses.
The homogeneous flat pivot Gaussian has the spectral and locality
estimates from Notebooks 6 and 7; Notebook 17 records its joint
fluctuation bounds. To apply the present theorem to the complete
Yang--Mills integrand, the action, Haar and Jacobian contributions
must be represented with the stated weighted local coefficient
norm on the actual cutoff domain.

The overlap, background-support, tunable-density and periodic
arguments below have passed separate internal mathematical audits.

